\chapter{Experiments and Evaluation}
\label{chap:experiments}

This chapter presents the experimental protocol used to evaluate the generated RTI fields. The objective is to compare generated samples with validation DNS fields using both generic feature-space metrics and physical indicators adapted to the Rayleigh--Taylor problem.

\section{Dataset Description}

The dataset consists of two-dimensional density fields extracted from Direct Numerical Simulations of Rayleigh--Taylor instability. The fields contain coherent large-scale structures and smaller-scale mixing patterns, while the number of available realizations remains limited because of the computational cost of DNS.

\textbf{TODO:} Add exact dataset details: number of simulations, number of snapshots, spatial resolution, train/validation/test split, physical parameters, mono-mode or bi-mode configuration, and storage format.

\section{Experimental Objectives}

The experiments compare the direct image-space diffusion model with the wavelet-based diffusion model. The evaluation focuses on visual realism, convergence behavior, statistical similarity to validation data, preservation of vertical density profiles, and reproduction of scale-dependent fluctuation spectra.

\section{Evaluation Metrics}

\section{Motivation}

The objective of the baseline diffusion model introduced in Chapter~3 is not to
reconstruct one specific DNS realization, but to generate new samples whose
distribution is close to that of the reference RTI dataset. Consequently, the
evaluation must be carried out at the level of \emph{ensembles} rather than on
individual images only.

In standard image generation tasks, this comparison is often performed with
feature-based distances such as the Fr\'echet Inception Distance (FID). Such
metrics are useful because they compare the statistics of two datasets in a
learned feature space. However, in the present context, the generated objects
are not natural images but scalar density fields arising from a physical
problem. It is therefore useful to complement generic perceptual metrics with
quantities that remain directly interpretable in terms of RTI structure.

For this reason, the evaluation adopted in this work combines two families of
metrics:

\begin{itemize}
    \item feature-based distribution distances, namely a classical
    Inception-based FID and a domain-adapted PhyFID;
    \item physically motivated one-dimensional summaries derived from the
    density fields, including line-wise mean profiles and fluctuation energy
    spectra.
\end{itemize}

This chapter presents the definition of these metrics and discusses the first
results obtained on the baseline DDPM.


\section{Feature-space metrics}

\subsection{Fr\'echet distance in a learned feature space}

Let $\rho=\{\rho_i\}_{i=1}^{n}$ be a reference dataset and
$\varrho=\{\varrho_i\}_{i=1}^{m}$ a generated dataset. A feature extractor
$\psi$ maps each image to a vector representation,

\begin{equation}
    \mathbf{z}_i = \psi(\rho_i),
    \qquad
    \widetilde{\mathbf{z}}_i = \psi(\varrho_i).
\end{equation}

The two sets of features are then approximated by Gaussian distributions with
means $\boldsymbol{\mu}_{\rho}$ and $\boldsymbol{\mu}_{\varrho}$, and
covariances $\mathbf{\Sigma}_{\rho}$ and $\mathbf{\Sigma}_{\varrho}$. The
Fr\'echet distance between these two Gaussian approximations is

\begin{equation}
    d_F(\rho,\varrho)
    =
    \left\|
    \boldsymbol{\mu}_{\rho}
    -
    \boldsymbol{\mu}_{\varrho}
    \right\|_2^2
    +
    \mathrm{Tr}\!\left(
        \mathbf{\Sigma}_{\rho}
        +
        \mathbf{\Sigma}_{\varrho}
        -
        2\left(\mathbf{\Sigma}_{\rho}\mathbf{\Sigma}_{\varrho}\right)^{1/2}
    \right).
\end{equation}

Low values indicate that the two datasets occupy nearby regions of feature
space and have similar second-order statistics. This metric does not compare
images pixel by pixel; instead, it compares the distributions of datasets after
projection onto a more compact representation.

\subsection{Classical FID and PhyFID}

Two choices of feature extractor are considered in this work.

The first one is the classical FID formulation, in which $\psi$ is an
Inception-style image encoder. This provides a generic image-generation
baseline and makes it possible to check whether the generated RTI fields are
distinguishable from completely unrelated datasets or from pure noise in a
standard visual feature space.

The second one is the \emph{PhyFID} introduced for this project. Instead of
using features learned on natural images, PhyFID relies on a convolutional
autoencoder trained directly on RTI fields. The encoder part of this network
defines the mapping $\psi$, while the decoder is used only during the
reconstruction training stage. After training, each dataset is encoded into a
feature space adapted to the statistics of the considered physical domain.

In practice, the PhyFID pipeline implemented in
\texttt{lib/PhyFID/} proceeds as follows:

\begin{itemize}
    \item the training images are normalized to the interval $[0,1]$ and
    converted to tensors of shape $(N,C,H,W)$;
    \item a convolutional autoencoder is trained by minimizing a reconstruction
    mean squared error;
    \item the encoder output is used as a domain-specific feature vector;
    \item the Fr\'echet distance is finally computed between the feature
    distributions of the reference and generated datasets.
\end{itemize}

This construction is especially relevant here because the characteristic
structures of RTI fields differ strongly from the objects encountered in
standard computer vision datasets. Using a feature extractor trained on the
same family of simulations should therefore provide a more meaningful notion of
similarity for this application.


\section{Physical metrics}

\subsection{Line-wise density statistics}

Chapter~3 highlighted the importance of the vertical organization of RTI
fields: the mixing region develops along the gravity direction, and the
position of the interface is one of the main physical signatures of the flow.
To preserve this information while reducing the dimensionality of the analysis, 
it is possible to compares the fields line by line (because of the horizontal periodicity).

Let $\rho(t,x,y)$ denote a validation field and $\varrho(t,x,y)$ a field from a
compared dataset, where $t$ indexes the images, $x$ the vertical coordinate and
$y$ the horizontal coordinate. For each image and each line, the mean density
along the periodic direction is first computed as

\begin{equation}
    m_{\rho}(t,x)
    =
    \frac{1}{N_y}
    \sum_{y=1}^{N_y}\rho(t,x,y),
    \qquad
    m_{\varrho}(t,x)
    =
    \frac{1}{N_y}
    \sum_{y=1}^{N_y}\varrho(t,x,y).
\end{equation}

The temporal mean profile of each line is then defined by

\begin{equation}
    \widetilde{\rho}_l(x)
    =
    \frac{1}{N_t}\sum_{t=1}^{N_t}m_{\rho}(t,x),
    \qquad
    \widetilde{\varrho}_l(x)
    =
    \frac{1}{N_t}\sum_{t=1}^{N_t}m_{\varrho}(t,x).
\end{equation}

This quantity represents the mean vertical density organization of the dataset.
It is directly related to the average position and thickness of the mixing
layer.

The temporal variability of each line is described by the standard deviation

\begin{equation}
    \sigma_{\rho,l}(x)
    =
    \sqrt{
        \frac{1}{N_t}
        \sum_{t=1}^{N_t}
        \left(
            m_{\rho}(t,x)-\widetilde{\rho}_l(x)
        \right)^2
    },
\end{equation}

and similarly for $\sigma_{\varrho,l}(x)$. This profile measures how strongly
the mean value of a given line fluctuates across the dataset.

The distances used in the notebook are normalized $L^2$ distances between these
profiles:

\begin{equation}
    d_{\mu}(\rho,\varrho)
    =
    \frac{1}{N_x}
    \left\|
        \widetilde{\rho}_l
        -
        \widetilde{\varrho}_l
    \right\|_2,
    \qquad
    d_{\sigma}(\rho,\varrho)
    =
    \frac{1}{N_x}
    \left\|
        \sigma_{\rho,l}
        -
        \sigma_{\varrho,l}
    \right\|_2.
\end{equation}

These metrics are simple but physically informative: they test whether the
generated ensemble reproduces the average vertical density profile and the
amplitude of its fluctuations.

\begin{figure}[H]
    \centering
    \includegraphics[width=0.90\textwidth]{figures/ch4/skewnorm_joyplot.png}
    \caption{Line-wise density statistics extracted from the validation,
    generated and noise datasets. The skew-normal fits emphasize the vertical
    organization of the fields and the contrast between physical samples and
    pure noise.}
    \label{fig:ch4_line_statistics}
\end{figure}

\subsection{Fluctuation 'energy' spectrum}

The term 'energy spectrum' is used here in a loose sense to describe the distribution of fluctuations across spatial scales. It is not a true physical energy spectrum, but it hase to be seen as the signale energy (it would have been equal to a real energy spectrum if the density field was a velocity field). The first physical metric introduced above captures the mean vertical organization of the fields, but it does not provide any information about the distribution of fluctuations across scales. Starting from a dataset converted to a tensor of shape
$(t,x,y)$, the fluctuation field is defined in the notebook by subtracting the
global mean and taking its magnitude,

\begin{equation}
    \delta \rho(t,x,y)
    =
    \left|
        \rho(t,x,y) - \overline{\rho}
    \right|.
\end{equation}

A discrete Fourier transform is then applied, and the energy spectrum is
computed as

\begin{equation}
    E_{\rho}
    =
    \left|
        \mathcal{F}(\delta\rho)
    \right|^2.
\end{equation}

In the baseline, a one-dimensional transform is applied along one
spatial direction before reusing the same line-wise statistical machinery as
above. The resulting quantities therefore compare the average spectral content
and its variability across the dataset. Even though this construction is not a
full isotropic turbulence spectrum, it provides a first indication of whether
the generated fluctuations occupy approximately the correct range of spatial
scales.

\begin{figure}[H]
    \centering
    \includegraphics[width=0.90\textwidth]{figures/ch4/skewnorm_joyplot_energy_spectrum.png}
    \caption{Same comparison applied to the fluctuation energy spectrum. The
    generated samples remain much closer to the validation set than pure noise,
    but the spectral variability is still harder to match than the mean line
    profile.}
    \label{fig:ch4_energy_spectrum}
\end{figure}


\section{Baseline results}

\subsection{PhyFID and FID comparison}

The notebook compares several datasets on 1000 images: a held-out part of the
validation set, a training subset, the generated DDPM samples, pure-noise
samples produced by an untrained model, and an unrelated MNIST dataset used as
a negative control. The obtained scores are reported in
Table~\ref{tab:fid_phyfid_baseline}.

\begin{table}[H]
    \centering
    \begin{tabular}{lcc}
        \hline
        Comparison & PhyFID & FID Inception \\
        \hline
        val vs val & 0.021 & 0.0013 \\
        val vs train & 0.023 & 0.0008 \\
        val vs generated & 0.132 & 0.134 \\
        val vs noise & 6.04 & 42.05 \\
        val vs MNIST & 4.26 & 14.27 \\
        \hline
    \end{tabular}
    \caption{Comparison of validation, training, generated, noise and MNIST
    datasets using PhyFID and classical FID on 1000 images. Lower values are
    better.}
    \label{tab:fid_phyfid_baseline}
\end{table}

Several conclusions can be drawn from these values. First, the two reference
comparisons `val vs val` and `val vs train` give very low distances. This is
expected and confirms that both metrics behave consistently when comparing
datasets drawn from the same RTI distribution.

Second, the generated dataset remains significantly closer to the validation
set than pure noise or MNIST. In PhyFID space, the score for `val vs generated`
is approximately two orders of magnitude smaller than `val vs noise`
($0.1315$ versus $6.0399$). The same qualitative trend is observed with the
classical FID, where `val vs generated` remains small compared with
`val vs noise` and `val vs MNIST`.

These results indicate that the baseline DDPM has learned non-trivial
statistical information about the RTI fields. The generated samples are clearly
not random images and it is therefore reasonable to continue the evaluation with 
other architectures and metrics.

\subsection{Line profile and spectral indicators}

The profiles and spectra extracted from the validation, generated and noise datasets are shown in Figures~\ref{fig:ch4_line_statistics} and
\ref{fig:ch4_energy_spectrum}. The corresponding distances are,

\begin{table}[H]
    \centering
    \begin{tabular}{lcc}
        \hline
        Comparison & mean & fluctuation energy spectrum \\
        \hline
        val vs val & 0 & 0 \\
        val vs generated & 0.49 & 0.298 \\
        val vs noise & 13.1 & 1.66 \\
        val vs MNIST & 10.7 & 0.735 \\
        \hline
    \end{tabular}
    \caption{Comparison of validation, generated, noise and MNIST datasets using line-wise mean profile and fluctuation energy spectrum distances on 1000 images. Lower values are better.}
    \label{tab:line_profile_spectrum_baseline}
\end{table}

The results confirm the conclusions drawn from the feature-space metrics. The generated samples are much closer to the validation set than pure noise or MNIST, both in terms of mean vertical profile and fluctuation energy spectrum with proximately two orders of magnitude.

\section{Baseline Results}

The baseline diffusion model generates samples that reproduce the main organization of the reference simulations and remain significantly closer to the validation dataset than pure noise or unrelated images. In PhyFID space, the generated fields obtain a distance of approximately .132$ from the validation set, compared with .04$ for pure noise.

\textbf{TODO:} Add final tables and figures for all reported metrics. Each figure and table must be explicitly cited and interpreted in the text.

\section{Wavelet-based Results}

The wavelet-based formulation introduces a multi-resolution inductive bias and shows that diffusion can learn relationships between large-scale approximations and fine-scale details. In its current implementation, reconstructed fields remain rougher than those obtained with the direct baseline, which indicates that the conditioning mechanism and architecture still need refinement.

\textbf{TODO:} Evaluate the wavelet-based model with the same metrics as the baseline: PhyFID, FID, line-wise statistics, and fluctuation spectra. Add qualitative comparisons of reference, baseline, wavelet reconstruction, and generated fields.

\section{Interpretation}

The current results show that direct diffusion is the strongest baseline in terms of generation quality, while wavelet diffusion provides a more interpretable and physically meaningful direction for future improvement. The comparison highlights the importance of adapting architectures and normalization strategies to transformed coefficient spaces.

This evaluation chapter showed the strengths and limitations of the implemented pipelines. The final chapter summarizes the main findings and discusses perspectives for improving the wavelet-based approach.
